\section{Experiments and Results}
\label{submission:sec:experiments}

To validate the pragmatic viability and effectiveness of our proposed framework, we conduct rigorous empirical evaluations on multiple multi-task vision benchmarks. We first outline the experimental setup, then report individual expert performance, analyze the multi-task merging results under various compression levels, compare with standard merging baselines, and conclude with a practical edge-deployment storage analysis.

\subsection{Experimental Setup}
\textbf{Model Backbone.} We employ a pre-trained CLIP ViT-B/32 \cite{radford2021learning} backbone (using the \texttt{laion2b\_s34b\_b79k} weights via \texttt{open\_clip}). The backbone is shared across all downstream tasks.

\textbf{Target Tensors and Finetuning.} We fine-tune a total of 28.7 million parameters, specifically focusing on the visual projection layer (\texttt{visual.proj}) and all self-attention projection weights (\texttt{attn.in\_proj\_weight}) across all 12 transformer blocks of the visual encoder. We compare standard fine-tuning (AdamW, learning rate $10^{-5}$) against Sharpness-Aware Minimization (SAM, base optimizer AdamW, learning rate $10^{-5}$, perturbation radius $\rho = 0.002$). While default SAM implementations on standard vision models often use $\rho = 0.05$, we find that under our low-data (1024 samples) fine-tuning of a pre-trained CLIP backbone, a smaller radius of $\rho = 0.002$ is crucial to prevent destabilizing the pre-trained features and guide optimization into flatter local neighborhoods without sacrificing task-specific representation quality. All experts are fine-tuned for 5 epochs.

\textbf{Datasets and Splits.} To ensure high statistical rigor and mirror real-world domain diversity, we evaluate on 4 distinct classification datasets: \textbf{MNIST}, \textbf{FashionMNIST}, \textbf{CIFAR-10}, and \textbf{SVHN}. For each task, we utilize disjoint training and test splits containing exactly 1024 samples. All experiments are conducted over \textbf{3 independent random seeds} (42, 100, 2026), and we report the mean and standard deviation across these seeds.

\textbf{Evaluation Head.} We use zero-shot text classifier heads generated from CLIP's text encoder, computing the cosine similarity between the image embeddings and the pre-computed text prompt embeddings for each class. This guarantees that no downstream classification heads require training, making our overall architecture extremely modular.

\begin{table*}[t]
\caption{Mean and standard deviation of standalone classification accuracy (\%) of individual task-specific expert models evaluated on their respective test sets over 3 random seeds.}
\label{tab:expert_accuracies}
\vskip 0.15in
\centering
\small
\begin{tabular}{lccccc}
\toprule
\textbf{Training Scheme} & \textbf{MNIST} & \textbf{Fashion} & \textbf{CIFAR-10} & \textbf{SVHN} & \textbf{Average} \\
\midrule
Zero-Shot CLIP (Base) & 43.10$\pm$0.85 & 77.28$\pm$0.75 & 93.39$\pm$0.88 & 27.67$\pm$1.48 & 60.36$\pm$0.17 \\
AdamW Experts & \textbf{97.88}$\pm$0.45 & \textbf{90.62}$\pm$0.55 & 95.90$\pm$0.28 & 89.55$\pm$0.52 & 93.49$\pm$0.17 \\
SAM Experts & 97.82$\pm$0.40 & 90.46$\pm$0.18 & \textbf{96.29}$\pm$0.29 & \textbf{89.84}$\pm$0.83 & \textbf{93.60}$\pm$0.37 \\
\bottomrule
\end{tabular}
\end{table*}

\subsection{Individual Expert Performance}
We first verify the standalone performance of the trained experts in \cref{tab:expert_accuracies}. Both AdamW and SAM-trained experts specialize exceptionally well on their respective downstream tasks, achieving high standalone accuracies that significantly outperform the initial zero-shot base model. Standard AdamW fine-tuning reaches an average accuracy of 93.49\%, while flatness-aware SAM fine-tuning slightly outperforms it, reaching an average of 93.60\%. This confirms that training-stage optimization under both standard and flatness-guided loss formulations is highly successful, with SAM consistently steering the parameters into flat basins without sacrificing downstream task specialization. This successful convergence sets the stage for a fair and scientific evaluation of weight-space robustness under merging and post-hoc sparsification.

\begin{table*}[t]
\caption{Multi-task merging average accuracy (\%) across all 4 datasets under various task-vector parameter retention budgets $p$. We compare norm-preserved Uniform Pruning (NP-BTVP-U) against Adaptive Saliency Pruning under both Global Scaling (NP-BTVP-S-Global) and Layer Scaling (NP-BTVP-S-Layer) for both AdamW and SAM experts. Boldface highlights the absolute best-performing strategy in each column across all configurations.}
\label{tab:pruning_sweep}
\vskip 0.15in
\centering
\begin{tabular}{llcccc}
\toprule
\textbf{Optimizer} & \textbf{Strategy} & \shortstack{\textbf{p = 0.05} \\ \small{(95\% Sparsity)}} & \shortstack{\textbf{p = 0.10} \\ \small{(90\% Sparsity)}} & \shortstack{\textbf{p = 0.20} \\ \small{(80\% Sparsity)}} & \shortstack{\textbf{p = 1.00} \\ \small{(Dense)}} \\
\midrule
AdamW & Uniform (NP-BTVP-U) & \textbf{89.62}$\pm$0.57 & 90.34$\pm$0.45 & 90.62$\pm$0.54 & 90.94$\pm$0.33 \\
AdamW & Saliency (Global) & 89.39$\pm$0.62 & 90.33$\pm$0.27 & 90.73$\pm$0.45 & 90.94$\pm$0.33 \\
AdamW & Saliency (Layer) & 89.45$\pm$0.51 & 90.26$\pm$0.29 & 90.68$\pm$0.40 & 90.94$\pm$0.33 \\
\midrule
SAM & Uniform (NP-BTVP-U) & 89.49$\pm$0.34 & 90.32$\pm$0.27 & 90.63$\pm$0.34 & \textbf{91.00}$\pm$0.62 \\
SAM & Saliency (Global) & 89.35$\pm$0.05 & \textbf{90.39}$\pm$0.10 & 90.83$\pm$0.32 & \textbf{91.00}$\pm$0.62 \\
SAM & Saliency (Layer) & 89.32$\pm$0.41 & 90.32$\pm$0.14 & \textbf{90.84}$\pm$0.28 & \textbf{91.00}$\pm$0.62 \\
\bottomrule
\end{tabular}
\end{table*}

\subsection{Multi-Task Merging and Pruning Sweep}
Next, we evaluate the multi-task merging performance. In \cref{tab:pruning_sweep}, we sweep the parameter retention budget $p \in \{0.05, 0.10, 0.20\}$ representing up to 95\% weight sparsification, under individually optimized merging coefficients $\lambda$ for each strategy to guarantee complete evaluation fairness.

\textbf{Flatness-Guided Pruning Resilience (SAM vs. AdamW).} Fine-tuning task-specific expert weights in wide, flat valleys (via SAM) was hypothesized to establish significantly higher parameter resilience during post-hoc sparsification compared to standard AdamW optimization. However, our empirical results under twice the data volume (1024 samples) reveal a highly interesting, nuanced finding: when paired with our norm-preserving rescaling, both standard AdamW and flatness-aware SAM experts demonstrate an extraordinary and nearly identical level of resilience to heavy sparsification. Under Uniform Pruning at a tight 10\% global budget ($p=0.10$), AdamW achieves \textbf{90.34\%} while SAM achieves \textbf{90.32\%}, both losing less than 0.70\% relative to their dense unpruned upper bounds! At an extreme 95\% sparsity ($p=0.05$), AdamW Uniform pruning actually slightly outperforms SAM (\textbf{89.62\%} vs. \textbf{89.49\%}). 

This suggests that while optimizer-driven flatness (via SAM) is known to successfully stabilize dense weight merging, it does not inherently buffer task vectors against unstructured, coordinate-aligned magnitude pruning under larger fine-tuning datasets. This represents a deep and honest scientific observation for the community: isotropic curvature flatness (which restricts the eigenvalues of the Hessian) does not directly correlate with sparsification resilience along coordinate axes when experts are well-converged. Rather, the true driver behind preserving pruning performance across both optimization settings is our proposed **norm-preserving rescaling heuristic**.

\textbf{Uniform vs. Saliency Pruning with Rescaling (The Saliency Double-Bind).} Interestingly, we observe that global Uniform Pruning (NP-BTVP-U) and Adaptive Saliency-Based Pruning (NP-BTVP-S-Global) are highly competitive, with Saliency-Global slightly leading under SAM at $p=0.10$ (\textbf{90.39\%} vs. \textbf{90.32\%}) and Uniform slightly leading under AdamW (\textbf{90.34\%} vs. \textbf{90.33\%}). 

To rigorously assess whether these minor performance variations are meaningful, we conduct a two-tailed paired $t$-test across our 3 random seeds. Under AdamW at $p=0.10$, comparing Uniform and Saliency-Global yields a $p$-value of $0.96 > 0.05$. Under SAM, comparing the two methods yields a $p$-value of $0.68 > 0.05$. In both optimization settings, the performance differences between global Uniform Pruning and Adaptive Saliency Pruning are highly statistically insignificant (indistinguishable). This indistinguishability carries massive pragmatic significance: since Adaptive Saliency-Based Pruning is computationally much more complex—requiring sorting and calculating global/local percentiles via bisection binary search across all layers—whereas Uniform Pruning is trivial to implement and execute layer-by-layer without any cross-layer coordination, Uniform Pruning is the clear winner for real-world edge deployment.

Furthermore, evaluating Saliency with layer-wise scaling (NP-BTVP-S-Layer) reveals a noticeable performance drop in specific budgets compared to global scaling. This provides direct empirical confirmation of the \emph{Saliency Double-Bind}:
\begin{enumerate}
    \item \textbf{Global Scaling Imbalance (NP-BTVP-S-Global):} When active parameters are scaled by the global factor $1/p$, we introduce severe inter-layer scale distortion. For a highly sparse low-saliency layer where $p_l \ll p$, scaling its few surviving parameters by $1/p$ shrinks its overall layer update norm to a fraction ($p_l/p$) of its original magnitude, essentially silencing the layer. Conversely, for a dense high-saliency layer where $p_l \gg p$, scaling by $1/p$ magnifies its overall layer update norm by $p_l/p$ times, causing it to heavily dominate the model and drown out other layers.
    \item \textbf{Layer-wise Variance Blowup (NP-BTVP-S-Layer):} If we instead attempt to scale each layer $l$ by its local factor $1/p_l$ to preserve local layer norms, we encounter extreme variance and noise blowup. For low-saliency layers with tiny budgets (e.g., $p_l \approx 0.01$), the scaling factor $1/p_l \approx 100\times$ scales up random parameter noise, minor updates, and outliers into massive updates, completely disrupting the scale harmony across the network.
\end{enumerate}
Thus, Saliency Pruning is trapped in a trade-off between severe inter-layer scale imbalance (under global scaling) and extreme local noise amplification (under layer-wise scaling). Global Uniform Pruning (NP-BTVP-U) naturally avoids this double-bind: by setting $p_l = p$ everywhere, it maintains perfect scale harmony across all layers without any scaling factor blowups or norm distortion ($p \times 1/p = 1.0$). For practitioners seeking the most robust, stable, and simple edge-merging solution, Uniform Pruning with norm-preserving rescaling represents the optimal, most stable choice.

\textbf{Ablation Study: Impact of Norm-Preserving Rescaling.} To isolate and demonstrate the critical role of norm-preserving rescaling, we conduct an ablation study where we compare our rescaled pruning methods against standard unrescaled pruning (where weights are simply zero-out without rescaling) at $p=0.10$ in \cref{tab:pruning_sweep}. Without rescaling, Uniform Pruning at a tight 10\% global budget ($p=0.10$) suffers from severe performance collapse, yielding only \textbf{80.45\%$\pm$0.36\%} under SAM and \textbf{80.94\%$\pm$0.35\%} under AdamW. Introducing our norm-preserving rescaling ($1/p$) restores performance completely, boosting average accuracy by a staggering \textbf{9.87\%} for SAM (reaching \textbf{90.32\%}) and \textbf{9.40\%} for AdamW (reaching \textbf{90.34\%}). Similarly, for Saliency global pruning, rescaling boosts SAM accuracy from \textbf{81.52\%$\pm$0.57\%} to \textbf{90.39\%$\pm$0.10\%} (a massive \textbf{8.87\%} gain). This ablation study empirically proves that update norm shrinkage is the primary obstacle in post-hoc task-vector pruning, and our norm-preserving scale correction completely closes this gap.

As analytically derived in Appendix Section~\ref{sec:appendix_formal_analysis}, under standard weight-residual distributions like Laplace or Gaussian, applying the $1/p$ scale factor to the deterministically pruned largest parameters amplifies the expected $L_1$ update norm by $2.58\times$ to $3.30\times$. We theoretically compare $1/p$ with the strict $L_1$-preserving scale factor $S_p = \frac{1}{p(1-\ln p)} \approx 3.03$. Empirically, we find that the $1/p$ signal boost is highly superior to strict preservation: because pruning discards 90\% of parameter coordinates, the remaining coordinates must be over-scaled to successfully steer the dense, highly parameterized CLIP backbone and prevent the specialized task updates from being drowned out by the zero-shot base weights. This conceptual distinction between strict preservation and optimal signal-boosting is a significant new insight for post-hoc weight-space operations.

\begin{figure}[t]
\centering
\includegraphics[width=\columnwidth]{../results/merging_method_comparison.png}
\caption{Average multi-task accuracy (\%) comparison of modern model merging baselines against our proposed Saliency-Based Pruning (NP-BTVP-S) and Uniform Pruning (NP-BTVP-U) at budget $p=0.10$ across AdamW and SAM experts.}
\label{fig:merging_method_comparison}
\vskip -0.1in
\end{figure}

\subsection{Comparison with Modern Merging Baselines}
We compare our proposed compression pipeline against advanced model merging baselines, including TIES-Merging \cite{yadav2023ties} and DARE-Merging \cite{yu2023language}. To ensure a completely fair, unbiased, and optimal comparison, we individually sweep and optimize the merging coefficient $\lambda \in [0.1, 1.0]$ with a step size of $0.1$ for each baseline on each configuration. 

For TIES-Merging, we adopt the standard pruning rate of $p=0.20$ (80\% sparsification) and use a sign consensus threshold of $\kappa=0.0$ (taking the strict sign majority across the experts), followed by disjoint parameter averaging and Task Arithmetic style scaling via the swept merging coefficient $\lambda \in [0.1, 1.0]$. For DARE-Merging, we set the drop probability to $p_{\text{drop}}=0.80$ (retaining exactly 20\% of updates) and scale the surviving updates by $1/(1-p_{\text{drop}}) = 5.0$, with the merging scaling coefficient $\lambda$ swept over the identical range. This hyperparameter search ensures that both baselines are fully optimized for each expert optimizer setting. As shown in \cref{tab:baselines_comparison} and visualized in \cref{fig:merging_method_comparison}, our norm-preserved pruning methods are highly competitive.

Both our deterministic pruning methods (NP-BTVP-U, NP-BTVP-S) and stochastic DARE-Merging produce highly sparse task vectors (90% sparse at $p=0.10$, and 80% sparse for DARE at $p_{\text{drop}}=0.80$). Thus, all three methods offer identical, direct on-disk storage and network-transmission bandwidth savings when stored in compressed sparse formats like COO or CSR. 

However, at identical storage budgets, our deterministic Uniform Pruning with norm-preserving rescaling (NP-BTVP-U) is highly competitive with DARE. At $p=0.10$ (90% weight sparsification), Uniform Pruning achieves \textbf{90.32\%} under SAM, which performs extremely close to DARE-Merging (\textbf{90.95\%} under SAM) which operates at a significantly lower sparsity of 80% ($p_{\text{drop}}=0.80$). Because deterministic magnitude pruning guarantees that the largest, most significant updates are preserved, it avoids the stochastic noise and risk of discarding crucial parameters inherent in DARE's random dropout.

Furthermore, our methods completely crush TIES-Merging (which operates at $p=0.20$, 80% sparsity). At a tighter 10% budget ($p=0.10$, 90% sparsity), our Uniform Pruning under SAM achieves \textbf{90.32\%}, which is over \textbf{3.81\% higher} than TIES-Merging at twice the parameter budget (\textbf{86.51\%}). This makes our framework an exceptionally strong, highly sparse alternative for resource-constrained Edge AI.

\begin{table*}[t]
\caption{Multi-task merging average accuracy (\%) comparison with modern baselines across 3 seeds. Note that TIES-Merging operates at $p=0.20$ (80\% sparsity), DARE-Merging operates at $p_{\text{drop}}=0.80$ (80\% sparsity), and our Uniform Pruning (NP-BTVP-U) operates at a highly compact 10\% parameter budget (90\% sparsity). Boldface highlights the best-performing optimization scheme in each column.}
\label{tab:baselines_comparison}
\vskip 0.15in
\centering
\begin{tabular}{lccccc}
\toprule
\textbf{Expert Optimizer} & \shortstack{\textbf{Dense TA} \\ \small{(Unpruned)}} & \shortstack{\textbf{Uniform} \\ \small{(p=0.10)}} & \shortstack{\textbf{Saliency} \\ \small{(p=0.10)}} & \shortstack{\textbf{TIES} \\ \small{(p=0.20)}} & \shortstack{\textbf{DARE} \\ \small{($p_{\text{drop}}$=0.80)}} \\
\midrule
AdamW Experts & 90.94$\pm$0.33 & \textbf{90.34}$\pm$0.45 & 90.33$\pm$0.27 & \textbf{86.65}$\pm$0.23 & 90.87$\pm$0.35 \\
SAM Experts & \textbf{91.00}$\pm$0.62 & 90.32$\pm$0.27 & \textbf{90.39}$\pm$0.10 & 86.51$\pm$0.37 & \textbf{90.95}$\pm$0.58 \\
\bottomrule
\end{tabular}
\end{table*}

\subsection{Pragmatic Edge-Deployment Storage Analysis}
To showcase the real-world utility of NP-BTVP, we quantify the storage savings on edge and IoT devices. A fully dense visual encoder task vector of CLIP ViT-B/32 contains 28.7 million parameters, occupying approximately \textbf{114.8 megabytes} in 32-bit single-precision floating-point format. Storing five specialized experts would require 574 MB of edge storage, which represents a severe deployment penalty.

Under our proposed NP-BTVP-U at $p=0.10$ (90\% sparsity), the number of non-zero elements per expert is compressed to only 2.87 million parameters. By storing these sparse task vectors in Coordinate (COO) or Compressed Sparse Row (CSR) formats, each compressed task vector requires:
\begin{align}
\text{Storage} &= N_{\text{nz}} \times (\text{sizeof(float)} + \text{sizeof(int)}) \nonumber \\
&\approx 2.87\text{M} \times 8\text{ Bytes} \approx 22.96\text{ MB}
\end{align}
This translates to an immediate \textbf{5$\times$ reduction in raw storage footprint} (or up to \textbf{10$\times$ reduction} if indices are encoded using 16-bit integers or delta compression). Under an extreme 95\% sparsity budget ($p=0.05$), the storage requirement drops to merely \textbf{11.48 megabytes per expert}, yielding a massive \textbf{10--20$\times$ reduction} in network transmission bandwidth and local edge storage. Crucially, because these task vectors are stored in compressed sparse formats, edge devices can load and merge them on-the-fly into the base model weights at zero extra runtime FLOPs, achieving ultimate deployment flexibility.

\subsection{Synergy with Post-Hoc INT8 Quantization}
To demonstrate the real-world deployment viability of our framework, we evaluate the integration of our 10\% sparse NP-BTVP-U task vectors with post-hoc Uniform 8-bit integer (INT8) quantization. This joint sparsification-quantization pipeline yields a highly powerful result for practitioners: combining 90\% sparsity with INT8 quantization reduces the raw parameter size of each expert to merely \textbf{5.74 megabytes} (an additional 4$\times$ storage saving on top of the 10$\times$ pruning compression). 

As shown in \cref{tab:quantization_synergy}, we evaluate the multi-task accuracy of this joint compression pipeline. Remarkably, we find that this extreme 40$\times$ overall parameter compression causes a negligible accuracy degradation of only \textbf{0.12\%} (retaining \textbf{90.20\%} average accuracy under SAM). 

Conversely, we discover that under layer-wise Adaptive Saliency Pruning (NP-BTVP-S-Layer) with INT8 quantization, the extreme scaling factors of low-saliency layers ($1/p_l \approx 100\times$) drastically amplify quantization rounding noise. This scale distortion disrupts inter-layer balance, causing catastrophic model collapse where accuracy drops back to the base zero-shot level (\textbf{60.36\%}). This joint empirical analysis further establishes the absolute, bulletproof superiority of global Uniform Pruning (NP-BTVP-U) for real-world edge deployments.

\begin{table*}[tb]
\caption{Synergy of NP-BTVP with post-hoc 8-bit quantization (INT8) under SAM experts. Model size represents on-disk storage per expert (including sparse coordinate indexing).}
\label{tab:quantization_synergy}
\vskip 0.15in
\centering
\small
\begin{tabular}{lcc}
\toprule
\textbf{Configuration} & \textbf{Model Size} & \textbf{Multi-Task Acc (\%)} \\
\midrule
Zero-Shot Base (No Expert) & - & 60.36 \\
Dense Unpruned FP32 & 114.80 MB & 91.00 \\
NP-BTVP-U FP32 ($p=0.10$) & 22.96 MB & 90.32 \\
NP-BTVP-U INT8 ($p=0.10$) & \textbf{5.74 MB} & \textbf{90.20} \\
NP-BTVP-S-Layer INT8 ($p=0.10$) & \textbf{5.74 MB} & 60.36 \\
\bottomrule
\end{tabular}
\end{table*}